	% -----------------------------------------------
% Styl pro psaní diplomových a bakalářských prací
% Určeno pro překlad pdfcsLaTeXem
% -----------------------------------------------
\documentclass[a4paper,12pt]{report} % pro oboustranný tisk zvolte {book}!
% velikost stránky
\usepackage[tmargin=2cm,bmargin=2.5cm,rmargin=2cm,lmargin=3.5cm]{geometry}
% volba kódování, nastaveno je unicode
\usepackage[utf8]{inputenc}
% \usepackage[cp1250]{inputenc}
% \usepackage[latin2]{inputenc}
\usepackage[english]{babel}
%\def\refname{Literatura}
% pomocná makra pro sazbu matematických výrazů
\usepackage{amssymb} %na psaní dvojitých písem a zvlastních znaků (např. \varkappa)
\usepackage{amsmath}
% balíček pro vkládání obrázků
\usepackage{graphicx}
% balíček pro obrázky na krajích stránky 
\usepackage{floatflt}
% balíček pro křížové odkazy
\usepackage[unicode,bookmarksopen,colorlinks=false,plainpages=false,urlcolor=blue,pdfpagelabels]{hyperref}
% balíček pro vkládání hypertextových odkazů
\usepackage{url}

\usepackage[numbers]{natbib}

% -----------------------------------------------
% Přepínač mezi bakalářskou a diplomovou prací, 
% barevným a černobílým logem a mezi studentem a studentkou (vypracoval/vypracovala apod.)
% -----------------------------------------------
\newif\ifbakal\bakaltrue
\newif\ifbarva\barvafalse
\newif\ifstudentka\studentkafalse

% -----------------------------------------------
% Údaje o práci – doplní se i do bibliografické identifikace, pouze v prohlášení je nutné vložit jméno vedoucího práce, které není v 1. pádě
% -----------------------------------------------
\newcommand{\nazevcz}{Kalibrace a monitorování astročásticových teleskopů}
\newcommand{\nazev}{Calibration and monitoring of astroparticle telescopes}
\newcommand{\student}{Daniel Staník}
\ifbakal%
  \newcommand{\program}{B0533A110007  Applied Physics}%
  \newcommand{\obor}{1702R001  Applied Physics  (AFYZ)}%
  %\else%
  %\newcommand{\program}{N1701 Fyzika}%
  %\newcommand{\obor}{7504T055 Učitelství fyziky pro střední školy}%
\fi
\newcommand{\rokod}{2022}
\newcommand{\vedouci}{Ing. Ladislav Chytka, Ph.D.}
\newcommand{\abstrakt}{%
Lorem ipsum dolor sit amet, consectetur adipiscing elit. Curabitur et lectus sit amet lectus vestibulum dignissim. Cras sit amet enim vitae mi elementum blandit eget nec tortor. Curabitur eget eros vitae arcu luctus varius commodo vel mauris. Nam elementum convallis pretium. Nunc dignissim pulvinar urna, nec blandit ante fringilla at. Ut et magna purus, vel pellentesque massa. In tortor nisi, faucibus condimentum cursus ut, sollicitudin quis leo. Ut at purus nec arcu accumsan tincidunt id id massa. Nam id vehicula mi.}
% -----------------------------------------------
\newcommand{\abstrakten}{%
Lorem ipsum dolor sit amet, consectetur adipiscing elit. Curabitur et lectus sit amet lectus vestibulum dignissim. Cras sit amet enim vitae mi elementum blandit eget nec tortor. Curabitur eget eros vitae arcu luctus varius commodo vel mauris. Nam elementum convallis pretium. Nunc dignissim pulvinar urna, nec blandit ante fringilla at. Ut et magna purus, vel pellentesque massa. In tortor nisi, faucibus condimentum cursus ut, sollicitudin quis leo. Ut at purus nec arcu accumsan tincidunt id id massa. Nam id vehicula mi.}
% klíčová slova
\newcommand{\klic}{klíčové slovo 1, klíčové slovo 2, \ldots}
\newcommand{\klicen}{keyword 1, keyword 2, \ldots}
\newcommand{\pocetstran}{xx}
\newcommand{\pocetpriloh}{x}

% -----------------------------------------------
% Definice vlastních maker pro usnadnění psaní
% a opakování symbolů při zlomu řádku
% -----------------------------------------------
% implikace se opakuje
\def\Plyne{\Rightarrow\discretionary{}{\hbox{$\Rightarrow$}}{}}
% -----------------------------------------------
% ekvivalence se opakuje
\def\Ekviv{\Leftrightarrow\discretionary{}{\hbox{$\Leftrightarrow$}}{}}
% -----------------------------------------------
% % plus '+' se opakuje při zalomení řádku
\mathchardef\plus="202B
\mathcode`\+="8000
{\catcode`\+=\active
\gdef+{\plus\nobreak\discretionary{}{\hbox{$\plus$}}{}}}
% opakování - při zalomení řádku
\newsavebox{\minusbox}
\savebox{\minusbox}{\hbox{$-$}}
\def\aktivniminus #1{{\catcode`#1=13 \bgroup \uccode`~=`#1
\uppercase{\egroup\gdef~}{\mathminus\discretionary{}{\copy\minusbox}{}}}}
% % -----------------------------------------------
% % rovnitko '=' se opakuje
\def\aktivnirovnitko #1{{\catcode`#1=13 \bgroup \uccode`~=`#1
\uppercase{\egroup\gdef~}{\mathequal\discretionary{}{=}{}}}}
% -----------------------------------------------
% zrušení mezery za čárkou v matematickém reľimu
\mathcode`,="002C

% -----------------------------------------------
% Úprava matematické sazby
% -----------------------------------------------
\def\sgn{\mathop{\rm sgn}\nolimits}
\def\tg{\mathop{\rm tg}\nolimits}
\def\cotg{\mathop{\rm cotg}\nolimits}
\def\arctg{\mathop{\rm arctg}\nolimits}
% -----------------------------------------------
% vektory polotučným skloněným písmem
\DeclareFontFamily{OT1}{bssbf}{}
\DeclareFontShape{OT1}{bssbf}{m}{n}{<5> <6> <7> <8> <9> <10> <11> <12> <14.4> <17> <20> <20.74> <25> bssb10}{}
\newcommand{\vek}{\fontencoding{OT1}\fontfamily{bssbf}\selectfont}
\renewcommand{\vec}[1]{\hbox{\vek #1}\hspace*{1.5pt}}
% polotučná řecká písmena
\def\bgomega{\vec{\char151}}
     \def\bgOmega{\vec{\char10}}
     \def\bggamma{\vec{\char130}}
     \def\bgvarphi{\vec{\char161}}
     \def\bgphi{\vec{\char148}}
     \def\bgxi{\vec{\char142}}
     \def\bgtau{\vec{\char146}}
     \def\bgeta{\vec{\char134}}
     \def\bgmu{\vec{\char139}}
     \def\bgnu{\vec{\char140}}
     \def\bvarrho{\vec{\char157}}
     \def\bgsigma{\vec{\char145}}
     \def\bgpsi{\vec{\char150}}
     \def\bgchi{\vec{\char149}}
     \def\bgvartheta{\vec{\char154}}

% ---------------------------------------------------------------
% Samotná práce
% ---------------------------------------------------------------
\begin{document}
% -----------------------------------------------
% Titulní strana
% -----------------------------------------------
\pagestyle{empty}
\setbox0=\hbox{\LARGE\scshape Joint Laboratory of Optics}
\centerline{\LARGE\scshape  Palacký University Olomouc}

\bigskip
\centerline{\LARGE\scshape Faculty of Science}
   
\bigskip
\centerline{\box0}
  
\vfill
\centerline{\LARGE\bfseries \ifbakal{BACHELOR}\else{DIPLOMOVÁ}\fi\ THESIS}

\bigskip
\begin{center}
{\huge\nazev}  


\vfill
\ifbarva\includegraphics[height=3cm]{up_logo_color}\else%
\includegraphics[height=3cm]{up_logo_bw}\fi

\vfill

\noindent%
\begin{tabular}{ll}
Author\ifstudentka{a}\fi: & {\bfseries\student}\\
Study program: &\program\\
Field of study: &\obor\\
Form of study:& Full-time\\
Supervisor:& \vedouci\\
Deadline:& April~\rokod
\end{tabular}
\end{center}

% ---------------------------------------------------------------
% Prohlášení
% -----------------------------------------------
\newpage
\hbox{~}

\vfill

\begin{center}
{\bf DECLARATION}
\end{center}

\noindent
I hereby declare that I elaborated this bachelor thesis independently under the supervision 
of  Ing. Ladislav Chytka, Ph.D.,  using  only  information  sources  referred  in  the  Literature chapter. \\
{}\vspace{3ex}

\noindent
In~Olomouc~\today   \hfill\parbox[t]{6cm}{\centering\null\dotfill\\\student}

% -----------------------------------------------
% Bibliografická anotace
% -----------------------------------------------
\newpage
% Bibliografická identifikace
\section*{Bibliographical identification}

\begin{tabular}{lp{8cm}}
%-----------
% \multicolumn{2}{c}{\bfseries Bibliographical identification}\\[8mm]
%-----------
Autor's first name and surname & \student\\
%-----------
Title & \nazev\\
%-----------
Type of thesis & \ifbakal{Bachelor}\else{Master}\fi \\
%-----------
Department & Joint Laboratory of Optics \\
%-----------
Supervisor & \vedouci\\
%-----------
The year of presentation & \rokod \\
%-----------
Abstract & \abstrakten\\
%-----------
Keywords & \klicen\\
%-----------
Number of pages & \pocetstran\\
%-----------
Number of appendices &  \pocetpriloh\\
%-----------
Language & czech\\
%-----------
\end{tabular}

% -----------------------------------------------

\newpage
\section*{Bibliografická identifikace}

\begin{tabular}{lp{8.5cm}}
%-----------
% \multicolumn{2}{c}{\bfseries Bibliografická identifikace}\\[8mm]
%-----------
Jméno a příjmení autora & \student\\
%-----------
Název práce & \nazevcz \\
%-----------
Typ práce & \ifbakal{Bakalářská}\else{Diplomová}\fi \\
%-----------
Pracoviště & Společná Laboratoř Optiky \\
%-----------
Vedoucí práce & \vedouci\\
%-----------
Rok obhajoby práce & \rokod\\
%-----------
Abstrakt & \abstrakt\\
%-----------
Klíčová slova & \klic\\
%-----------
Počet stran & \pocetstran\\
%-----------
Počet příloh & \pocetpriloh\\
%-----------
Jazyk & český\\
%-----------
\end{tabular}
% %%%%%%%%%%%%%%%%%%%%%%%% End of file %%%%%%%%%%%%%%%%%%%%%%%%





%%%% Tady začínáme %%%%%%%%%%%%%%%%%%%%%%%%%%%%%%%%%%%%%%%%%%%%%%%%%%%%%%%%%%
\newpage
%%%%%%%%%%%%%%%%%%%%%%%%%%%%%%%%%%%%%%%%%%%%%%%%%%%%%%%%%%%%%%%%%%%%%%%%%%%%%
\widowpenalty =10000
\pagestyle{plain}
% -----------------------------------------------
% Obsah je generován automaticky, změny se projeví po 2 překladech
% -----------------------------------------------
\tableofcontents

% -----------------------------------------------
% Úvod
% -----------------------------------------------
\chapter*{Introduction}
\addcontentsline{toc}{chapter}{Introduction}

Lorem ipsum dolor sit amet, consectetur adipiscing elit. Curabitur et lectus sit amet lectus vestibulum dignissim. Cras sit amet enim vitae mi elementum blandit eget nec tortor. Curabitur eget eros vitae arcu luctus varius commodo vel mauris. Nam elementum convallis pretium. Nunc dignissim pulvinar urna, nec blandit ante fringilla at. Ut et magna purus, vel pellentesque massa. In tortor nisi, faucibus condimentum cursus ut, sollicitudin quis leo. Ut at purus nec arcu accumsan tincidunt id id massa. Nam id vehicula mi. 

\url{http://exfyz.upol.cz/didaktika/}

% -----------------------------------------------
% Kapitoly lze ukládat do zvláštních souborů...
% -----------------------------------------------
% -----------------------------------------------
% Vlastní text práce (kapitoly práce)
% -----------------------------------------------

% -----------------------------------------------
\chapter{Astroparticle detection}
% -----------------------------------------------
More than 100 years have passed since Victor Franz Hess first encontered cosmic radiation. Since those times the techniques and methods of detection have been strongly improved. We have moved up from elevating electroscopes by ballons to observe growing electric charge to specialized techniques, which allows us to measure particles' energies, trajectories, etc.

% -----------------------------------------------
\section{Cosmic rays and particles}
% -----------------------------------------------
Cosmic rays is a term for radiation and energetic particles striking earth atmosphere with an origin in a space sources (neutron stars, supernovas, black holes, etc). We divide them into two major groups - primary and secondary. Primary cosmic rays are the original cosmic particles, which strike the Earth's athmosphere. Secondary cosmic rays (also refered as showers) are particles, which have origin in particle interaction between primary cosmic rays and the athmosphere.
\par
\subsection{Primary cosmic rays}
Primary cosmic rays consist of protons ($95 \%$), helium nuclei ($4 \%$), electrons and other heavy nuclei (up to iron). However, only the energetic rays make their way to the athmosphere. The Earth's magnetic field affects their trajectories and prevents the low-energetic (less than 100 MeV) particles from arriving to the athmosphere \cite{Kliewer}. 
\par
Part of primary cosmic rays are also Ultra-high energy cosmic rays (UHECRs), which we refer in the next chapter.
\par 
Neutrinos are also a part of cosmic radiation, but their interaction with matter is very rare, so they are very hard to detect. The special underwater detectors are developed to detect some of them. 
\subsection{Secondary cosmic rays}
Secondary cosmic rays are created by interaction of high-energetic particles of primary component with air's nucleis, such as nitrogen. They consist of low-energetic and high-energetic muons, gamma photons, electrons and positrons. Most of muons travel up to the earth's surface although their half-life is only about 2.2 microseconds before they decay into electrons. Due to their high relativistic speeds, their half-life is increased for external observers. 




% -----------------------------------------------
\section{Ultra-high energy cosmic rays (UHECRs)}
UHECRs are particles with energies from $10^{18}$ to $10^{20}$ eV, which is much more than particles created on Cern's Large hadron colider (LHC) with energies about $10^{13}$. Due to their high energies, the trajectory remains nearly unchanged by space magnetic fields \cite{Benjamin_Skuse}.
\par
UHECRs' origin is yet unknown, but it is supposed and experimentally proved that they come from outside of the Milky Way. Some theoretical physicists expects, that the one possible source of UHECRs acceleration are the starburst galaxies.
\par
When the UHECR interacts with athmospheric nuclei,

% -----------------------------------------------
\section{UHECRs detection techniques}


\subsection{Pierre Auger observatory}
% -----------------------------------------------

\subsection{Telescope array project}
% -----------------------------------------------
% %%%%%%%%%%%%%%%%%%%%%%%% End of file %%%%%%%%%%%%%%%%%%%%%%%%

% -----------------------------------------------
% -----------------------------------------------
% Vlastní text práce (kapitoly práce)
% -----------------------------------------------

% -----------------------------------------------
\chapter{FAST telescope}
% -----------------------------------------------
The  Fluorescence  detector  Array  of  Single-pixel  Telescopes  (FAST) is an international project of fluorescence telescope sensitive to UHERCs. 
\par
Until today there are four prototypes in active service. Three of them are situated in Black Rock Mesa site of the Telescope Array experiment in central Utah and one in Argentina near Pierre Auger Observatory.
\par
The main goal of FAST project is to develop a cheap fluorescence telescope, which could be used in future to cover the wide surface area. This new oncoming fluorescence telescope array should be able  
to fully reconstruct the geometry of UHERCs induced UV shover by combining the information from telescopes, which has encountered some event at the same time. 
% -----------------------------------------------
\section{Principle of operation}
% -----------------------------------------------
Main detection part of telescope consists of superreflective UV mirrors and photomultipliers. 


\par
The entire telescope along with monitoring systems and other instruments is situated in a hut with remote shutter, where it is protected from negative metrological phenomena, such as rain or fast wind, but also from dust and aerosols. Exposure of mirrors to any of this phenomena could lead to reduction of theirs reflectivity. It is also neccessary to monitor and protect PMTs from unwanted light sources. Even a low-intensity sources could decrease PMT's service life.
% -----------------------------------------------

\section{Remote control and monitoring}
% -----------------------------------------------

%------------------------------------------------

\section{}
% -----------------------------------------------



% -----------------------------------------------
% %%%%%%%%%%%%%%%%%%%%%%%% End of file %%%%%%%%%%%%%%%%%%%%%%%%

% -----------------------------------------------
% Závěr
% -----------------------------------------------
\chapter*{Conclusion}
\addcontentsline{toc}{chapter}{Conclusion}

Lorem ipsum dolor sit amet, consectetur adipiscing elit. Curabitur et lectus sit amet lectus vestibulum dignissim. Cras sit amet enim vitae mi elementum blandit eget nec tortor. Curabitur eget eros vitae arcu luctus varius commodo vel mauris. Nam elementum convallis pretium. Nunc dignissim pulvinar urna, nec blandit ante fringilla at. Ut et magna purus, vel pellentesque massa. In tortor nisi, faucibus condimentum cursus ut, sollicitudin quis leo. Ut at purus nec arcu accumsan tincidunt id id massa. Nam id vehicula mi. 

% -----------------------------------------------
% Literatura a prameny
% -----------------------------------------------
%\begin{thebibliography}{99}
\bibliographystyle{apsrev4-2}
\nocite{*}
\bibliography{citations}
%\addcontentsline{toc}{chapter}{Literatura}
%\bibitem{gravitation} MISNER, Ch. W.; THORNE, K. S.; WHEELER, J. A. %\emph{Gravitation}. San Francisco: W. Freeman, 1973.
%\end{thebibliography}

\newpage
Preferované jsou citace podle norem ČSN ISO 690 a ISO 690-2, popř. styly APS (American Physical Society – u~prací zaměřených fyzikálně) nebo APA (American Psychological Association – u~prací zaměřených více didakticky a pedagogicky).
\end{document}
% Konec souboru %%%%%%%%%%%%%%%%%%%%%%%%%%%%%%%%%%%%%%%%%%%%%%%%%%%%%%%%%%%%
